\section{Reproducing this report}
\label{app:repro}

\begin{table}[htbp]
\centering
\caption{Everything in this document is generated. The order below is the order
to run it in. The fitting step dominates; the rest takes under two minutes.}
\label{tab:repro}
\small
\begin{tabular}{@{}p{5.4cm}p{9.4cm}@{}}
\toprule
\textbf{Command} & \textbf{What it produces} \\
\midrule
\code{pip install -r requirements-dev.txt} & The full analysis environment.
\code{requirements.txt} alone is the deployment runtime \\[3pt]
\code{python tests/test\_leakage.py} & 41 validation checks: join grain, the
history construction against a brute-force recomputation, the leakage guard in
both directions, split disjointness, Phase 1 backward compatibility, and the
capacity and lateness tables against recomputation \\[3pt]
\code{python src/train\_models.py} & Fits every model, scores the test window
once, writes \code{models/} and every result table into \code{report/results/} \\[3pt]
\code{python src/capacity\_metrics.py} & Precision@$K$ and Recall@$K$ from the
saved per-order test predictions, for Section~\ref{sec:results-capacity};
nothing is refitted \\[3pt]
\code{python src/lateness\_anatomy.py} & The weekly, monthly and strike-cohort
lateness tables of Appendix~\ref{app:lateness} \\[3pt]
\code{python src/figures\_phase2.py} & Draws every figure as vector PDF into
\code{report/figures/} from those tables \\[3pt]
\code{python src/phase2\_stats.py} & Distils every number quoted in this report
into \code{report/stats\_phase2.txt}. Diff it against the manuscript \\[3pt]
\code{python src/report\_stats.py} & Still reproduces the \emph{Phase 1} numbers
unchanged --- the check that the Phase 2 feature additions are backward
compatible \\[3pt]
\code{streamlit run dashboard/app.py} & The dashboard, including the Predict tab
that serves the models above \\
\bottomrule
\end{tabular}
\end{table}

\code{SEED = 42} is set in every estimator, every randomised search and every
subsample. Both requirements files pin exact versions on Python $\le$ 3.12, the
interpreter every reported number was produced on. Model artifacts record the
\code{scikit-learn} version that wrote them and \code{modeling.load\_model()}
checks it on the way back in.

\section{Success criteria registered in advance}
\label{app:bands}

Written into \code{config.EXPECTED\_BANDS} before any Phase 2 model was tuned.
\code{evaluate.check\_bands()} compares the headline results against them at the
end of every training run and writes its verdicts to
\code{report/results/bands.csv}. The alarm is direction-aware: for an error
metric the suspicious end of the interval is the low one, for a score the high
one, and confusing the two would point the check at the shortfalls and wave the
leaks through.

\begin{center}
\small
\begin{tabular}{@{}lllp{6.3cm}@{}}
\toprule
\textbf{Problem} & \textbf{Metric} & \textbf{Registered band} & \textbf{Where the range came from} \\
\midrule
1A & MAE (days) & 3.3 -- 4.3 & Phase 1 reported 3.89 untuned \\
1B & balanced accuracy & 0.58 -- 0.72 & Phase 1 reported 0.58 for a direct classifier \\
1B & ROC-AUC & 0.72 -- 0.85 & Untuned scoping probe reached 0.79 for the derived route \\
2 & balanced accuracy & 0.63 -- 0.74 & Phase 1 reported 0.69 \\
2 & ROC-AUC & 0.70 -- 0.79 & Phase 1 reported 0.73 \\
2 & PR-AUC & 0.34 -- 0.48 & Phase 1 reported 0.40 against a 0.11 floor \\
\bottomrule
\end{tabular}
\end{center}

\section{Hyperparameter search spaces and selected values}
\label{app:hyperparameters}

\begin{center}
\footnotesize
\begin{tabular}{@{}llll@{}}
\toprule
\textbf{Family} & \textbf{Rung} & \textbf{Classification} & \textbf{Regression} \\
\midrule
A --- Linear     & 0 & \code{LogisticRegression} & \code{LinearRegression} \\
                 & 1 & \code{LogisticRegression} ($\ell_2$, tuned $C$) & \code{Ridge} (tuned $\alpha$) \\
B --- Tree-based & 0 & \code{DecisionTreeClassifier} & \code{DecisionTreeRegressor} \\
                 & 1 & \code{RandomForestClassifier} & \code{RandomForestRegressor} \\
                 & 2 & \code{HistGradientBoostingClassifier} & \code{HistGradientBoostingRegressor} \\
C --- Ensemble   & 0 & \code{VotingClassifier} (soft, A+B) & \code{VotingRegressor} (A+B) \\
                 & 1 & \code{StackingClassifier} (A+B) & \code{StackingRegressor} (A+B) \\
\bottomrule
\end{tabular}
\end{center}

Every tuned model was searched with \code{RandomizedSearchCV} over the
distributions below, scored by five-fold expanding-window cross-validation inside
the training period. Thirty draws per model, except the random forests, which get
six: a forest costs one to two orders of magnitude more per fit than the other
rungs, and Phase 1 had already shown the forest and the boosted trees landing
within 0.01 days of each other. That is a documented trade rather than a hidden
one.

\begin{table}[htbp]
\centering
\caption{Search spaces. Selected values are in
\code{report/results/reg\_p1a.csv}, \code{clf\_p1b.csv} and \code{clf\_p2.csv},
and are reproduced in \code{models/model\_metadata.json}.}
\label{tab:spaces}
\small
\begin{tabular}{@{}llp{7.1cm}@{}}
\toprule
\textbf{Model} & \textbf{Draws} & \textbf{Distributions} \\
\midrule
\code{LinearRegression} / \code{LogisticRegression} & --- &
No free parameters. Fitted and scored under the same protocol so the row is
comparable \\[3pt]
\code{Ridge} & 30 & $\alpha \sim \mathrm{loguniform}(10^{-3}, 10^{3})$ \\[3pt]
\code{LogisticRegression} ($\ell_2$) & 30 & $C \sim \mathrm{loguniform}(10^{-3}, 10^{2})$ \\[3pt]
\code{DecisionTree} & 30 & \code{max\_depth} $\sim$ randint(3, 20);
\code{min\_samples\_leaf} $\sim$ randint(5, 200) [classifier: (20, 400)];
\code{max\_features} $\sim U(0.4, 1.0)$ \\[3pt]
\code{RandomForest} & 6 & 150 trees; \code{max\_depth} $\sim$ randint(8, 18)
[classifier: (6, 16)]; \code{min\_samples\_leaf} $\sim$ randint(20, 140)
[classifier: (30, 180)]; \code{max\_features} $\sim U(0.15, 0.40)$ \\[3pt]
\code{HistGradientBoosting} & 30 & \code{max\_iter} $\sim$ randint(120, 500);
\code{learning\_rate} $\sim \mathrm{loguniform}(0.02, 0.3)$;
\code{max\_leaf\_nodes} $\sim$ randint(15, 64);
\code{min\_samples\_leaf} $\sim$ randint(10, 120);
\code{l2\_regularization} $\sim \mathrm{loguniform}(10^{-4}, 10)$ \\[3pt]
\code{Voting} & --- & Weights fixed at 1:2 in the tree's favour. A stated prior,
not a tuned parameter: under drift the linear member is materially weaker and an
equal blend would be worse than its better half \\[3pt]
\code{Stacking} & --- & Meta-learner \code{Ridge}($\alpha=1$) /
\code{LogisticRegression}; inner CV five contiguous folds --- see the note below \\
\bottomrule
\end{tabular}
\end{table}

\paragraph{A compromise inside stacking, stated plainly.}
\code{StackingClassifier} and \code{StackingRegressor} build their meta-features
with \code{cross\_val\_predict}, which requires the inner folds to \emph{partition}
the training set: every row must be predicted exactly once. An expanding window
partitions nothing --- the earliest rows are never in a test fold --- so the
closest valid approximation is unshuffled \code{KFold}, which at least keeps each
fold contiguous in time. The meta-learner is consequently fitted on out-of-fold
predictions partly made by models trained on later orders.

The residual optimism is confined to the meta-learner's coefficients; the base
models are refitted on the whole training period afterwards, and the outer
protocol scoring the finished ensemble is still the expanding window. It matters
less than it otherwise would because of how the result lands: the ensemble is
\emph{flattered} by this arrangement and still does not win, which makes the
conclusion in Section~\ref{sec:results-ensemble} safer rather than weaker.

\section{Re-checking a Phase 1 claim about least squares}
\label{app:encoding}

Phase 1 reported that ordinary least squares was unusable on this design matrix
--- ``rank-deficient, so its coefficients explode''. Phase 2 re-ran it four ways.
The observation reproduces exactly; the conclusion drawn from it needed
correcting.

\begin{table}[htbp]
\centering
\caption{Reproducible via \code{report/results/encoding\_diagnostic.csv}. The
deficiency is exactly four --- one per categorical, because each one-hot block
sums to one and is therefore collinear with the intercept.}
\label{tab:encoding}
\small
\begin{tabular}{@{}lccrrr@{}}
\toprule
\textbf{Encoding} & \textbf{rank / cols} & \textbf{deficiency} &
\textbf{max $|\beta|$} & \textbf{CV MAE} & \textbf{test MAE} \\
\midrule
OLS, full one-hot block   & 101 / 105 & 4 & $2.7\times10^{13}$ & $3.9\times10^{11}$ & 4.181 \\
OLS, one level dropped    & 101 / 101 & 0 & 14.7  & 5.242 & 4.181 \\
Ridge, full one-hot block & 101 / 105 & 4 & 9.82  & 5.241 & 4.182 \\
Ridge, one level dropped  & 101 / 101 & 0 & 12.5  & 5.240 & 4.183 \\
\bottomrule
\end{tabular}
\end{table}

Least squares still returns a minimum-norm solution and its \emph{predictions}
are fine --- 4.181 days on the test window, indistinguishable from Ridge. What is
not fine is the coefficient vector at $10^{13}$, and the cancellation holding it
together is fragile enough to fail on one of five expanding-window folds, which
is why the cross-validated error is astronomical while the single-split test
error looks healthy. A single train/test split would never have shown it.

Dropping one level per categorical restores full rank, so the linear family is
encoded that way. The compromise is small and worth stating: an all-zero row is
also how the encoder represents a level it has never seen, so under
\code{drop="first"} a novel payment type is scored as the reference level. One
payment level is unseen in the test window, affecting a handful of orders.
Leaving the design rank-deficient costs more.

The last row carries the real lesson, and it is not the one Phase 1 drew. Once
the design is full rank, Ridge and OLS agree to three decimals. Regularisation
buys numerical stability here, not accuracy.

\section{Full result tables}
\label{app:tables}

The body reports the models that matter. The complete tables --- every model,
every metric, cross-validated, validation and test scores side by side --- are
written by \code{src/train\_models.py} to \code{report/results/}:

\begin{center}
\small
\begin{tabular}{@{}lp{7.8cm}@{}}
\toprule
\textbf{File} & \textbf{Contents} \\
\midrule
\code{reg\_p1a.csv} & Problem 1A: baselines, both ladders, ensembles; MAE, RMSE, $R^2$ \\
\code{clf\_p1b.csv} & Problem 1B: five direct classifiers and two derived routes \\
\code{clf\_p2.csv} & Problem 2: both ladders and both ensembles, all metrics \\
\code{encoding\_diagnostic.csv} & The Phase 1 re-check of Table~\ref{tab:encoding} \\
\code{cv\_protocol.csv} & Temporal against shuffled CV against the forward test \\
\code{weight\_ablation.csv} & Class weighting: discrimination and calibration \\
\code{calibration\_p2.csv}, \code{threshold\_sweep\_p2.csv} & Reliability and the operating-point sweep \\
\code{cost\_p2.csv}, \code{cost\_p1b.csv} & Expected net benefit by threshold \\
\code{perm\_p1a.csv}, \code{perm\_p2.csv} & Permutation importance on the test window \\
\code{segments\_*.csv} & Performance by region, category and test month \\
\code{predictions\_*.csv} & Per-order test predictions, from which the figures are drawn \\
\code{bands.csv} & Headline metrics against the bands registered in advance \\
\code{capacity\_*.csv} & Precision@$K$ and Recall@$K$, pooled and per day, for both flags; \code{capacity\_operating\_points.csv} places the threshold operating points on the same axis \\
\code{lateness\_*.csv} & The lateness anatomy of Appendix~\ref{app:lateness}: monthly, weekly, the strike cohorts and the February--March 2018 peak \\
\bottomrule
\end{tabular}
\end{center}

\section{The cross-validation protocol exhibit}
\label{app:protocol}

\begin{table}[htbp]
\centering
\caption{One model, three protocols. Raw average precision disagrees; the same
numbers divided by each period's no-skill floor do not.
Figure~\ref{fig:protocol}.}
\label{tab:protocol}
\small
\begin{tabular}{@{}lrrrr@{}}
\toprule
\textbf{Protocol} & \textbf{PR-AUC} & \textbf{fold spread} & \textbf{no-skill floor} & \textbf{lift} \\
\midrule
\code{TimeSeriesSplit}, expanding window & 0.562 & $\pm$0.091 & 0.158 & 3.55$\times$ \\
\code{StratifiedGroupKFold}, shuffled    & 0.580 & $\pm$0.006 & 0.158 & 3.66$\times$ \\
Forward test window --- the truth        & 0.414 & ---        & 0.110 & 3.75$\times$ \\
\bottomrule
\end{tabular}
\end{table}

Both cross-validated estimates sit about 0.16 above the forward result, which
looks like catastrophic optimism --- and almost all of it is mechanical. Average
precision's no-skill baseline \emph{is} the positive rate
\citep{saito2015prplot}, and the low-review rate falls from 15.8\% inside the
training period to 11.0\% in the test window; normalised by that floor the three
agree to within 0.2 of a multiple. \textbf{Raw PR-AUC is not comparable across
periods with different base rates; its lift over the floor is} --- and quoting
the raw number alone would have supported either a claim of severe overfitting or
a claim of none. The shuffled protocol adds a second warning: its fold-to-fold
spread is $\pm$0.006 against the temporal protocol's $\pm$0.091, an interval
fifteen times narrower on an estimate no closer to the truth.

\section{Lateness anatomy: where the late deliveries came from}
\label{app:lateness}

Phase 1 read the monthly late-rate series and attributed its spikes to Black
Friday and to the May 2018 truckers' strike. The Milestone 1 feedback pointed
out that the figure does not support the second claim -- the largest peaks are
February and March 2018, and May is unremarkable -- and asked for a source and a
more careful link. \code{src/lateness\_anatomy.py} is that analysis, at three
resolutions, and Figure~\ref{fig:lateness} draws it.

\paragraph{Weekly, not monthly.} A ten-day event cannot be resolved by a
monthly series. In Monday-to-Sunday purchase weeks the strike of 21 May to
1 June 2018 \citep{bloomberg2018strike} -- nationwide road blockades, during
which Correios suspended its time-definite Sedex services and added five
business days to standard parcels \citep{tecnoblog2018correios} -- is plainly
visible, but one week \emph{earlier} than a naive reading expects: orders placed
on 14--20 May were promised 19.5 days under normal conditions and were in
transit when the roads closed, and 16.0\% of them arrived late against a 5.0\%
average over the preceding five weeks. Of all orders in transit during the
strike (4{,}228), 16.9\% were late, against 0.4\% for the May orders delivered
before it (Table~\ref{tab:strike}). What the monthly series records instead is
the platform's reaction: order volume in the strike week halved, the promised
window for orders placed that week rose to 25.3 days and in the following week
to 36.8, and the late rate of those orders was 5.5\% and 0.4\%. That is why June
2018, at 1.4\%, is the quietest month in the data.

\paragraph{February--March 2018 is a different mechanism.} Panel (b) of
Figure~\ref{fig:lateness} splits lead time at the carrier handoff. Outside the
peak, sellers hand parcels to the carrier in about 3.2 days and the carrier
takes 8.7 more; in February--March 2018 the seller's share was unchanged at
3.4 days while the carrier's rose to 13.2, in every customer region. At the
same time the promised window was \emph{tightened}, from 26.3 days in January
to 22.0 in March -- within the S\~{a}o Paulo to S\~{a}o Paulo route alone, from
20.1 to 15.2 days -- so the slack that normally absorbs a slow fortnight was
removed exactly when it was needed. The late rate peaked at 29.0\% in the week
of 26 February, the promise was lengthened again from the week of 9 April, and
by that week the late rate was back under 5\%. The data do not identify the
external cause of the carrier slowdown and we do not assert one; what they do
establish is that it was carrier-side, network-wide and compounded by a
promising decision, which is the careful form of the claim Phase 1 made
loosely.

\begin{table}[htbp]
\centering
\caption{The cohorts around the strike, from
\code{report/results/lateness\_strike.csv}. Carrier days run from the carrier
handoff to the doorstep; the promise is the window quoted at checkout.}
\label{tab:strike}
\small
\begin{tabular}{@{}lrrrr@{}}
\toprule
\textbf{Cohort} & \textbf{orders} & \textbf{late} & \textbf{carrier days} & \textbf{promise, days} \\
\midrule
Bought 1--20 May, delivered before the strike & 2{,}749 & 0.4\%  & 4.8  & 19.5 \\
In transit during the strike                  & 4{,}228 & 16.9\% & 13.1 & 25.5 \\
Bought during the strike, 21 May--1 June      & 1{,}502 & 3.5\%  & 9.2  & 29.4 \\
Bought 2--30 June                             & 5{,}912 & 1.4\%  & 6.7  & 27.5 \\
\bottomrule
\end{tabular}
\end{table}

\section{Changes made in response to the Milestone 1 feedback}
\label{app:feedback}

\begin{table}[htbp]
\centering
\caption{Each point of the Milestone 1 feedback, and where this report acts on
it.}
\label{tab:feedback}
\small
\begin{tabular}{@{}p{4.4cm}p{10.1cm}@{}}
\toprule
\textbf{Feedback} & \textbf{Response} \\
\midrule
Survey-based perceived quality and platform ratings were treated as one construct &
Section~\ref{sec:problem} separates them, cites the self-selection literature
\citep{hu2009jshaped, hu2017selfselection, dellarocas2003digitization}, and
defines the target as the platform rating; Section~\ref{sec:limitations}
carries the consequence \\[2pt]
Model choices cited original algorithm papers rather than e-commerce
applications; no comparison within the tree family; HistGradientBoosting not
justified &
Section~\ref{sec:method} argues the rungs from the empirical comparisons and
from e-commerce applications, including three studies on this dataset, and
states why HistGradientBoosting; the ladder results in Section~\ref{sec:results}
are the within-family comparison \\[2pt]
Features not supported by literature &
Section~\ref{sec:data} cites the support for each feature group and uses
Phase 1's correlation matrix as the empirical screen \\[2pt]
Problem 2 should be evaluated under limited capacity, with Precision@$K$ and
Recall@$K$ &
Section~\ref{sec:results-capacity}, Table~\ref{tab:capacity} and
Figure~\ref{fig:capacity}; the operating point in Section~\ref{sec:cost} is now
a capacity; the dashboard's batch scorer returns a top-$K$ triage list \\[2pt]
Figure 3 attributed the late-rate peaks to the May 2018 strike; the strike
needed a source &
Appendix~\ref{app:lateness} and Figure~\ref{fig:lateness}: the peaks are
February--March 2018; the strike is cited \citep{bloomberg2018strike,
tecnoblog2018correios}, shown at weekly resolution, and its effect traced
through the promised window; Sections~\ref{sec:problem} and
\ref{sec:limitations} corrected accordingly \\[2pt]
Figure 5a omitted S\~{a}o Paulo; ``roughly triples''; Figure 8 not referenced &
Corrected in the revised Phase 1 report in the repository: the figure had taken
the first eight states alphabetically; growth is about eight-fold; the
correlation matrix is now discussed there and reused in Section~\ref{sec:data} \\[2pt]
Candidate 1 named a stakeholder without saying how the team would use the
prediction &
Table~\ref{tab:levers} and Section~\ref{sec:insight}: four levers ordered by
cost, each with its trigger and the evidence behind it \\[2pt]
Dashboard inaccessible &
Redeployed from the team's new repository under a new URL (cover page);
Appendix~\ref{app:deployment} explains how it stays reachable \\
\bottomrule
\end{tabular}
\end{table}

\section{Deployment}
\label{app:deployment}

The dashboard at \url{https://it5006-group-project-10.streamlit.app} carries the
five Phase 1 exploratory tabs unchanged and adds two.

\textbf{Predict} scores a single order entered as it would look at checkout ---
route, basket, timing and payment --- and returns the expected lead time, the
risk of missing the promised date against the tuned threshold, and the resulting
low-review risk. It also scores an uploaded CSV and returns it with the three
predictions appended. Nothing there reimplements the training-time feature
engineering: a form submission is turned back into a one-row order frame and
pushed through the same \code{features.py} builders the models were fitted on,
which is what stops a feature being computed one way in training and another way
in production. The one thing a form cannot supply is a seller's track record, so
\code{models/history\_lookup.csv} carries those aggregates at each level of the
same fallback chain \code{add\_history()} uses. After the Milestone 1 feedback
the batch scorer also returns a \textbf{triage list}: the uploaded orders ranked
by review risk, the top-$K$ marked for a capacity the user enters, and the
precision and recall that capacity achieved on the test window shown beside it
(Section~\ref{sec:results-capacity}).

\textbf{Model card} reports the training window, the selected models and their
hyperparameters, the test metrics, the decision thresholds, the Precision@$K$
table, and an explicit list of situations in which the models should not be
trusted.

\textbf{Staying reachable.} The Milestone 1 dashboard became inaccessible when
the GitHub account it was deployed from was suspended; the app now deploys from
the team's new repository, linked on the cover. A Community Cloud app that has
had no visitors for twelve hours goes to sleep and shows a wake-up button; any
visitor can press it, and the app is back in about a minute.

Engineering constraints are recorded rather than discovered later. The deployed
estimators are the histogram-boosted variants, whose pickles are a few hundred
kilobytes; a random forest of the same accuracy pickles to 24--30\,MB, which is
hostile to a shared container. Total artifact size is under 2\,MB.
\code{requirements.txt} pins \code{scikit-learn} to the exact version the models
were trained on and the loader raises a visible warning if it ever drifts. Models
are loaded lazily, so a visitor who only browses the exploratory tabs never pays
for unpickling one. Invalid input --- a promised date before the purchase date, a
negative price, more sellers than items, a malformed or empty upload, a missing
column --- produces a specific message rather than a traceback. No credentials
appear anywhere in the repository.

\section{Contributions}
\label{app:contributions}

\begin{center}
\small
\begin{tabular}{@{}lp{11.2cm}@{}}
\toprule
\textbf{Member} & \textbf{Contribution} \\
\midrule
Pan Pinyou & Problem scoping and stakeholder framing; feature engineering for the
track-record features; business interpretation and the cost analysis; report
Sections~\ref{sec:problem} and \ref{sec:insight} \\[3pt]
Li Yudan & Modelling pipeline and the family ladders; hyperparameter search;
imbalance and calibration analysis; report Sections~\ref{sec:method} and
\ref{sec:results-review} \\[3pt]
Zhu Gengyu & Data pipeline and the leakage guard; cross-validation protocol and
the direct-against-derived comparison; dashboard, deployment and reproducibility
infrastructure; report Sections~\ref{sec:data}, \ref{sec:results-reg} and
\ref{sec:limitations} \\
\bottomrule
\end{tabular}
\end{center}

All three members reviewed every section and verified the reported numbers
against \code{report/stats\_phase2.txt}.
